
\section*{Methodology}

\subsection*{Carbon Nanotube (CNT)}

The discovery of carbon nanotubes by Ijima of NEC Japan in 1993 is one potential
substitute for non-silicon electronics \cite{ref12}. When compared to steel,
they have higher tensile strength, electrical conductivity, thermal
conductivity, and nearly 1D ballistic transport capabilities
\cite{ref11,ref30,ref31,ref32,ref33,ref34}. Additionally, they have higher
electrical and thermal conductivity than diamond and steel. The chiral vector
defines the chirality, or helicity, of carbon nanotubes, and the atomic
composition can be described by the chiral angle ($\theta$). Figure 1 (a) shows
a graphene sheet with a known chiral vector and angle. Single-wall carbon
nanotubes (SWCNT) and multi-wall carbon nanotubes (MWCNT) are the two types of
carbon nanotubes shown in Figure 1 (b). However, Eq. (1) may demonstrate the
wrapping of graphite sheets in CNT, while Eq. (2) provides the lattice
translational indices (n, m) and CNT diameter (DCNT). The chiral vector, also
known as the roll-up vector, can be expressed in terms of the $a$1, $a$2 as unit
vectors. As seen in Figure 1(b), they can have two main forms: single-walled
CNTs (SWCNTs) and multi-walled CNTs (MWCNTs). MWCNTs are made up of several
graphene cylinders, whereas SWCNTs are made up of just one. Because of their
remarkable strength and unique electrical properties, CNTs can be used in a
variety of transdisciplinary domains, including electronics, optics, and
nanotechnology.

\begin{equation}
  Ch = n$a$1+m$a$2                                                                                                 (1)
  DCNT = Ch / π
  \label{eq:2}
\end{equation}

Three categories of single-walled carbon nanotubes (SWCNTs) are distinguished by
the chiral indices (n,m), which are integers that denote the steps along the
hexagonal lattice. Armchair nanotubes arise when n = m (with n \textgreater 0),
whose chiral vectors lie precisely between a1 or a2, corresponding to a chiral
angle of $\theta$ = 30∘. Zigzag nanotubes are produced when one of the chiral
indices, n or m, is zero. This results in a chiral vector aligned strictly along
either a1 or a2 direction, which corresponds to a chiral angle of $\theta$ = 0∘.
The remaining nanotubes are categorized as chiral, with chiral angles ranging
between 0∘ and 30∘, where n $\neq$ m \cite{ref6}.

\begin{figure}[htbp]
  \centering
  \begin{subfigure}[b]{0.48\linewidth}
    \centering
    \zimg{rf_fig_0.jpg}{width=\linewidth,max height=0.4\textheight,keepaspectratio}{fg_0_qwfkr}{rf_fig_0.jpg}
    \caption{(a) A representation of an unrolled graphene sheet along with the chiral vector definition (b) SWNT and MWNT}

  \end{subfigure}
  \hfill
  \begin{subfigure}[b]{0.48\linewidth}
    \centering
    \zimg{rf_fig_1.jpg}{width=\linewidth,max height=0.4\textheight,keepaspectratio}{fg_1_6r7t0}{rf_fig_1.jpg}
    \caption{(a) A representation of an unrolled graphene sheet along with the chiral vector definition (b) SWNT and MWNT}

  \end{subfigure}
  \caption{(a) A representation of an unrolled graphene sheet along with the chiral vector definition (b) SWNT}

  \label{fig:group_3e5ag}
\end{figure}

The CNT diameter, the inter-nanotube spacing (S), the number of CNTs in a
CNFET's channel (N), and the transistor's width (W) are all related by Equation
(3). Eq. (4) shows that the unit electron charge (e), carbon-to-carbon bond
energy (t; 3.0 eV), and carbon-to-carbon bond distance (acc; 0.1412 nm for
graphite) are all related to the band gap energy ($\sum$g) and threshold voltage
(Vth) of the intrinsic CNT channel \cite{ref7}.

\begin{equation}
  W = (N – 1)*S + DCNT                                                                                                            (3)
  ∑g = 2acct / DCNT(nm) = 0.84eV / DCNT(nm) ≈ 2eVth
  \label{eq:4}
\end{equation}

\subsection*{Carbon Nanotube field-effect transistor (CNTFET)}

One of the most important applications of carbon nanotubes is in carbon nanotube
field effect transistors (CNTFET). Figures 1(c) and 1(d) depict a CNTFET's
three-dimensional structure and top perspective, respectively, along with a few
device parameters. With its gate, source, drain, and substrate, the CNTFET is a
four-terminal device that resembles a MOSFET. While the source and drain
portions of CNTFETs are extensively doped, similar to MOSFETs, the CNT channel
region stays undoped, acting as both source/drain extension regions and/or
connectors between two neighboring devices \cite{ref17}. The gate controls the
channel's activation and deactivation behaviors. The MOSFET and CNFET exhibit
very similar operational behaviors. Like MOSFETs, CNFETs rely on their gate
terminals to apply a field perpendicular to the charge flow between their source
and drain terminals, which controls the carrier concentration in the channel
\cite{Imranetal2012} \cite{ref32}. Charge carriers travel ballistically across
CNTs from the drain to the source. A mean free path longer than the device's
size is indicated by ballistic transport. As a result, charge carriers do not
clash, which greatly lowers resistance to insignificant levels and increases
mobility in comparison to MOSFETs \cite{ref18}. Unlike CMOS technology, CNFET
technology uses a "1" device ratio for both p-type and n-type devices to design
circuits since it can drive the same amount of current using the same
geometries. The beneficial characteristics of CNTFETs include one-dimensional
ballistic charge carrier movement, high mobility, high drive current, and low
power consumption \cite{ref7,ref18,ref35}.

\begin{figure}[htbp]
  \centering
  \begin{subfigure}[b]{0.48\linewidth}
    \centering
    \zimg{rf_fig_2.png}{width=\linewidth,max height=0.4\textheight,keepaspectratio}{fg_0_13l9b}{rf_fig_2.png}
    \caption{(c) CNTFET's three-dimensional structure (d)The CNTFET's top view with its parameters}

  \end{subfigure}
  \hfill
  \begin{subfigure}[b]{0.48\linewidth}
    \centering
    \zimg{rf_fig_3.png}{width=\linewidth,max height=0.4\textheight,keepaspectratio}{fg_1_jw0k5}{rf_fig_3.png}
    \caption{(c) CNTFET's three-dimensional structure (d)The CNTFET's top view with its parameters}

  \end{subfigure}
  \caption{(c) CNTFET's three-dimensional structure (d)The CNTFET's top view with its parameters}

  \label{fig:group_zhbfs}
\end{figure}

\begin{itemize}
  \item Using field effect transistors based on carbon nanotubes, an operational
  transconductance amplifier: The operational transconductance amplifier (OTA),
  which is basically a voltage control current source, is a crucial component in
  the signal processing chain. The operational transconductance amplifier (OTA)
  can be defined as an amplifier with low impedance at every node except the
  input and output nodes \cite{ref1,ref2}. The cascode arrangement of an OTA
  offers high input impedance, high gain, and good stability. Reverse
  transmission, also known as input-output isolation, is enhanced since there is
  no direct coupling from the output to the input. The primary advantages of an
  OTA are its tunability and flexibility, which are really made possible by an
  external bias current (Ibias). Because of its versatility, an OTA can be used
  for a wide range of unique applications that aren't possible with the
  traditional OP-AMP, increasing its potential applications. Signal processing
  systems can be developed with OTAs \cite{ref5}. This functionality compensates
  for age, parasitics, fabrication tolerances, and variations caused by varying
  power supply and temperatures. Electrically tunable oscillators, filters,
  amplifiers, and current-mode signal processing devices can be made with OTAs
  \cite{ref24}. An ideal OTA's output resistance is infinite, and its output
  current (Iout) can be computed using equation (5) \cite{ref2}.
\end{itemize}

\begin{equation}
  Iout = gm (V1 – V2)
  \label{eq:5}
\end{equation}

where V1 and V2 are the non-inverting and inverting input terminals,
respectively, and gm is the transconductance \cite{ref6}.
Despite their many benefits, Operational Transconductance Amplifiers (OTAs) have
inherent drawbacks, especially when device dimensions approach the submicron
regime. Reduced channel length in typical OTAs causes a noticeable decline in
speed and gain, mostly because intrinsic gain (gmro) is significantly reduced.
The overall performance of the amplifier is negatively impacted by this
degradation \cite{ref4,ref7}. One commonly used method to address this problem
is cascading. Voltage gain and output impedance are significantly increased by a
cascode architecture, which is created by cascading a common-source (CS) stage
with a common-gate (CG) stage. When paired with a CG stage, the cascade
efficiently reduces the gain loss seen in conventional OTA designs. In a CS
setup, the transistor transforms the input voltage into an output current. In
the upper and lower branches of the suggested architecture, extra cascode
transistors are purposefully added to increase gain even more. Voltage gain and
overall amplifier performance are greatly enhanced by increasing the number of
cascoded devices \cite{ref28}. Specifically, a theoretical voltage gain of
around (gmro)2/2 is attained by the Triple Cascode Operational Transconductance
Amplifier (TC-OTA) architecture shown in Figure 2. The output current (Iout)
primarily passes through the external load capacitor (CL) due to its incredibly
high output impedance node (ROUT=RocasnǁRocasp), The TC-OTA design is
well-suited for high-gain, high-speed analog applications because of this
feature, which guarantees improved high-frequency performance and effective
signal transfer \cite{ref36}.

\begin{figure}[htbp]
  \centering
  \begin{subfigure}[b]{0.23\linewidth}
    \centering
    \zimg{rf_fig_4.png}{width=\linewidth,max height=0.3\textheight,keepaspectratio}{fg_0_bamv0}{rf_fig_4.png}
    \caption{(a) Pure CNTFET TC-OTA (b) NCNTFET-PMOS-TC-OTA (c) PCNTFET-NMOS-TC-OTA (d) Bulk CMOS-based TC-OTA}

  \end{subfigure}
  \hfill
  \begin{subfigure}[b]{0.23\linewidth}
    \centering
    \zimg{rf_fig_5.png}{width=\linewidth,max height=0.3\textheight,keepaspectratio}{fg_1_cv2lo}{rf_fig_5.png}
    \caption{(a) Pure CNTFET TC-OTA (b) NCNTFET-PMOS-TC-OTA (c) PCNTFET-NMOS-TC-OTA (d) Bulk CMOS-based TC-OTA}

  \end{subfigure}
  \hfill
  \begin{subfigure}[b]{0.23\linewidth}
    \centering
    \zimg{rf_fig_6.png}{width=\linewidth,max height=0.3\textheight,keepaspectratio}{fg_2_mhma5}{rf_fig_6.png}
    \caption{(a) Pure CNTFET TC-OTA (b) NCNTFET-PMOS-TC-OTA (c) PCNTFET-NMOS-TC-OTA (d) Bulk CMOS-based TC-OTA}

  \end{subfigure}
  \hfill
  \begin{subfigure}[b]{0.23\linewidth}
    \centering
    \zimg{rf_fig_7.png}{width=\linewidth,max height=0.3\textheight,keepaspectratio}{fg_3_dath8}{rf_fig_7.png}
    \caption{(a) Pure CNTFET TC-OTA (b) NCNTFET-PMOS-TC-OTA (c) PCNTFET-NMOS-TC-OTA (d) Bulk CMOS-based TC-OTA}

  \end{subfigure}
  \caption{(a) Pure CNTFET TC-OTA (b) NCNTFET-PMOS-TC-OTA (c) PCNTFET-NMOS-TC-OTA (d) Bulk CMOS-based TC-OTA}

  \label{fig:group_r0i95}
\end{figure}

The circuits of the three suggested CNT-based TC-OTAs and a conventional CMOS
triple cascode OTA are displayed in Figure 2. A pure CNT-TC-OTA is depicted in
Figure 2(a). It uses pCNTFETs as sources and nCNTFETs as sinks. A hybrid TC-OTA
known as NCNTFET-PMOS-TC-OTA is depicted in Figure 2(b) and uses nCNTFETS as
sinks and traditional PMOS transistors as sources. Similar to this, figure 2(c)
depicts another hybrid TC-OTA called PCNTFET-NMOS-TC-OTA that uses pCNTFETS as
sources and conventional NMOS transistors as sinks. As seen in figure 2(d), the
conventional CMOS based TC-OTA is made up completely of conventional NMOS and
conventional PMOS transistors. Every TC-OTA has been simulated using HSPICE and
uses a 32 nm technology node at 0.9V. In actuality, the simulation is
calibrated. As stated in \cite{ref16}, we used the experimental data of a
manufactured CNTFET to calibrate our CNTFET model. The structure, size, and
other characteristics of our simulated device are maintained in the model
calibration, just as they were used in [Javey et al. 2004] to generate
simulation data \cite{ref16}. The CNTFET reported in [Javey et al. 2004] employs
Pd source/drain contacts and an Al gate contact, with an intrinsic device
capacitance of approximately 2–5 aF per nanotube, along with overlap and fringe
capacitances of about 0.1 fF per nanotube. In this work, the Stanford University
CNTFET model is used for simulation. This SPICE-compatible compact model
accurately represents enhancement-mode, unipolar MOSFET-like devices utilizing
single-walled carbon nanotubes as the channel material. The model is based on
quasi-ballistic transport and incorporates several practical non-ideal effects,
including charge-carrier scattering due to acoustic and optical phonons,
parasitic capacitances, gate-to-gate and gate-to-contact plug capacitances,
charge-screening effects among adjacent nanotubes, source/drain resistances, and
band-to-band leakage currents \cite{ref7,ref16,ref26}.
The key simulation parameters include a physical channel length (Lch) of 45 nm,
an effective mean free path in the CNT (Lgeff) of 200 nm, a Fermi level of 0.6
eV for the doped source/drain nanotubes, a gate-oxide dielectric constant of 16,
a gate-oxide thickness (TOX) of 4 nm, a coupling capacitance of 40 pF·m$\^{-}$¹,
and a CNT work function of 4.5 eV. These parameters are consistent with those
employed in earlier CNT-based OTA designs \cite{ref7,ref15}. Since n-type and
p-type CNTFETs of identical dimensions exhibit comparable current-driving
capabilities, a unity sizing ratio between p-CNTFETs and n-CNTFETs is adopted in
the design of the TC-OTAs.

\subsection*{Impact of altering the number of CNTs (N) on the proposed TC-OTAs' performance}

It has been investigated how the performance of the proposed cascode OTAs is
affected by changes in the number of CNTs (N). N has been found to be crucial in
maximizing the effectiveness of the suggested TC-OTAs. Figure 4 illustrates how
N affects dc gain, output resistance (Rout), average power, and bandwidth. As
the number of CNTs increases, the dc gain in the CNT-based TC-OTAs first rises
and then saturates, as seen in Figure 3(a). The increase in driving capability
of both the PCNTFETs and the NCNTFETs is responsible for the initial rise in dc
gain. According to equation (6), CNTFETs' driving capability greatly increases
as N rises \cite{ref7,ref24,ref25}.

\begin{equation}
  ICNTFET≈NgCNT(VDD-Vth)/(1+gCNTLSρS)
  \label{eq:6}
\end{equation}

The dc gain becomes saturated as N increases due to the growth in screening
effect between the many CNTs in a CNTFET's channel, where $\rho$S is the source
resistance per unit length of doped CNT and Ls is the source length (doped CNT
region), gCNT is the transconductance per CNT, and ICNTFET is the CNTFET’s ON
current \cite{ref13,ref14,ref15,ref16,ref17,ref18,ref30}. Because a CNTFET's
effective channel width (Weff), as stated in equation (7) \cite{ref33},
decreases with N, its driving capability diminishes as N grows. Consequently, dc
gain decreases.

\begin{equation}
  Weff = N*DCNT* α
  \label{eq:7}
\end{equation}

where the screening effect coefficient $\alpha$ ranges from zero to one. The
screening effect coefficient $\alpha$ is equal to 1 when N = 1 and less than 1
when N \textgreater 1. Among all TC-OTAs, pure CNT-TC-OTA has a considerable
advantage over hybrid TC-OTAs due to the use of NCNTFETs and PCNTFETs, which
have superior transconductance over conventional bulk MOSFETs. Because it
comprises more CNTFETs, the NCNTFET-PMOS-TC-OTA has a bigger transconductance
than the hybrid PCNTFET-NMOS-TC-OTA, which explains why it has a higher gain in
the hybrid TC-OTAs. Equation (8) provides the formula for dc gain.

\begin{equation}
  DC gain (AV) = gmn (ROUT) = gmn (RocasnǁRocasp)
  \label{eq:8}
\end{equation}

where the transconductance of bulk NMOS transistors is lower than the gmn
transconductance of NCNTFET. Consequently, the highest dc gain is found in pure
CNT-TC-OTA; however, NCNTFET-NMOS-TC-OTA has a higher dc gain than
PCNTFET-NMOS-TC-OTA. Rocasn and Rocasp, respectively, represent the output
resistance of cascode-connected n-type and p-type devices \cite{ref1}.

\begin{figure}[htbp]
  \centering
  \begin{subfigure}[b]{0.48\linewidth}
    \centering
    \zimg{rf_fig_8.png}{width=\linewidth,max height=0.4\textheight,keepaspectratio}{fg_0_hmh6q}{rf_fig_8.png}
    \caption{(b) shows that the output resistance in all TC-OTAs diminishes as N grows. Since an increase in N is equivalent to an increase...}

  \end{subfigure}
  \hfill
  \begin{subfigure}[b]{0.48\linewidth}
    \centering
    \zimg{rf_fig_9.png}{width=\linewidth,max height=0.4\textheight,keepaspectratio}{fg_1_z31ws}{rf_fig_9.png}
    \caption{(b) shows that the output resistance in all TC-OTAs diminishes as N grows. Since an increase in N is equivalent to an increase...}

  \end{subfigure}
  \caption{(b) shows that the output resistance in all TC-OTAs diminishes as N grows. Since an increase in N is equivalent to an increase in a transistor's width W, output resistance (ROUT) must drop as N increases. However, TC-OTAs have a greater ROUT due to the cascoding effect. Figure 3(c) shows how N affects the average power consumption of CNT-based TC-OTAs. It is evident that average power grows along with N. This is explained by the fact that CNTFETs' driving capacity has increased due to the increase in N. However, because to the exceptionally low CNT channel resistance, 1D ballistic transport, reduced parasitic capacitance, and reduced leakage, a pure CNT-based TC-OTA has a much lower power dissipation \cite{ref18,ref35}. When comparing the two hybrid TC-OTAs, the NCNTFET-PMOS-TC-OTA uses more power than the PCNTFET-NMOS-TC-OTA because the former has more nCNTFETs (14 nCNTFETs) than the latter (12 pCNTFETs). The bandwidth is greatly enhanced in CNT-based TC-OTAs, as shown in Figure 3(d), with the amount of CNTs having the biggest impact on pure CNT-TC-OTA. It is more important for PCNTFET-NMOS-TC-OTA circuits than NCNTFET-PMOS-TC-OTA circuits, though, and grows linearly \cite{ref6,ref19}.}

  \label{fig:group_ab5t3}
\end{figure}

\begin{figure}[htbp]
  \centering
  \begin{subfigure}[b]{0.48\linewidth}
    \centering
    \zimg{rf_fig_10.png}{width=\linewidth,max height=0.4\textheight,keepaspectratio}{fg_0_wlkgp}{rf_fig_10.png}
    \caption{Impact of changing the number of CNTs in the proposed CNT-based TC-OTAs on (a) DC gain (b) output resistance, (c) average power, and (d) bandwidth}

  \end{subfigure}
  \hfill
  \begin{subfigure}[b]{0.48\linewidth}
    \centering
    \zimg{rf_fig_11.png}{width=\linewidth,max height=0.4\textheight,keepaspectratio}{fg_1_jkwal}{rf_fig_11.png}
    \caption{Impact of changing the number of CNTs in the proposed CNT-based TC-OTAs on (a) DC gain (b) output resistance, (c) average power, and (d) bandwidth}

  \end{subfigure}
  \caption{Impact of changing the number of CNTs in the proposed CNT-based TC-OTAs on (a) DC gain (b) output resistance, (c) average power, and (d) bandwidth}

  \label{fig:group_us7eu}
\end{figure}

\subsection*{Impact of altering the number of CNT diameter (DCNT) on the proposed TC-OTAs' performance}

The proposed pure CNT-TC-OTA and hybrid TC-OTAs can perform even better by
fine-tuning the CNT diameter used in various CNTFETs. Performance evaluations of
the suggested CNT-based TC-OTAs have been conducted for N = 20, and all of them
have been compared. Figure 4(a) shows a considerable variation in the dc gain of
the pure CNT-TC-OTA with DCNT. The first dc gain augmentation is caused by
strong transconductance and lower screening and scattering effects
\cite{ref6,ref30}. However, as DCNT gets bigger (\textgreater 1.4 nm), the gain
quickly drops due to an increase in screening and scattering effects. Figure
4(b) shows that in NCNTFET-PMOS-TC-OTA and PCNTFET-NMOS-TC-OTA, the output
resistance remains constant as the CNT diameter increases. In contrast, it rises
sharply for DCNT (1.0-1.2 nm) in pure CNT-TC-OTA, falls for DCNT (1.2-1.4 nm),
and stays constant for DCNT \textgreater 1.4 nm. This might be because smaller
CNTs have higher ROUT and fewer screening and scattering effects
\cite{ref6,ref13,ref14,ref15,ref16,ref17}. Additionally, the effect of ROUT
reduction outweighs the rise in transconductance as the CNT diameter increases.
The average power dissipation in CNT-based TC-OTAs and its variation with DCNT
are shown in Figure 4(c). As DCNT increases, both the NCNTFET-NMOS-TC-OTA and
PCNTFET-NMOS-TC-OTA circuits show a significant rise in power dissipation.
Moreover, an increase in DCNT initially causes a linear decrease in power
dissipation in pure CNT-TC-OTA, according to recent observations. But for DCNT
\textgreater 1.2 nm, the power dissipation significantly increases. Likewise,
the increased driving capability is due to the increase in transconductance.
Furthermore, an increase in CNT diameter reduces the bandgap in accordance with
Equation (8) \cite{ref14,ref18,ref35,ref36,ref37}.

\begin{equation}
  Bandgap = 0.8 eV / DCNT
  \label{eq:8}
\end{equation}

In CNTFETs, a lower threshold voltage is caused by a smaller bandgap, which
improves leakage and static power dissipation. Furthermore, the average power of
NCNTFET-PMOS-TC-OTA rises more than that of PCNTFET-NMOS-TC-OTA as the CNT
diameter increases. The increase in NCNTFET-PMOS-TC-OTA driving current brought
about by the inclusion of additional CNTFETs (14 NCNTFETs) justifies this. The
bandwidth of CNT-based TC-OTAs increases with increasing CNT diameter, as shown
in Figure 4(d). This can be explained by the fact that ROUT decreased due to the
increasing DCNT, as shown by Equation (9). Furthermore, increasing the CNT
diameter results in a rise in transconductance, which expands the bandwidth of
TC-OTAs based on CNTs \cite{ref16}.

\begin{equation}
  f3dB = 1/2πROUTCL
  \label{eq:9}
\end{equation}

The bandwidth gain in pure CNT-TC-OTA is more noticeable than in
PCNTFET-NMOS-TC-OTA and NCNTFET-PMOS-TC-OTA. Since the bandgap and, thus, the
threshold voltage of the CNTFET are large, the driving current will decrease for
CNTs with a lower diameter. The modest driving current in pure CNT-TC-OTA
results in a limited bandwidth for small CNT diameters. Such a huge CNT diameter
will improve the driving current and, eventually, the bandwidth in pure
CNT-TC-OTA by reducing the bandgap and, consequently, the threshold voltage.

\begin{figure}[htbp]
  \centering
  \begin{subfigure}[b]{0.48\linewidth}
    \centering
    \zimg{rf_fig_12.png}{width=\linewidth,max height=0.4\textheight,keepaspectratio}{fg_0_ilyjg}{rf_fig_12.png}
    \caption{Impact of changing the CNT diameter in the proposed CNT-based TC-OTAs on (a) DC gain (b) output resistance, (c) average power, and (d) bandwidth}

  \end{subfigure}
  \hfill
  \begin{subfigure}[b]{0.48\linewidth}
    \centering
    \zimg{rf_fig_13.png}{width=\linewidth,max height=0.4\textheight,keepaspectratio}{fg_1_ibi0z}{rf_fig_13.png}
    \caption{Impact of changing the CNT diameter in the proposed CNT-based TC-OTAs on (a) DC gain (b) output resistance, (c) average power, and (d) bandwidth}

  \end{subfigure}
  \caption{Impact of changing the CNT diameter in the proposed CNT-based TC-OTAs on (a) DC gain (b) output resistance, (c) average power, and (d) bandwidth}

  \label{fig:group_uoudr}
\end{figure}

\begin{figure}[htbp]
  \centering
  \begin{subfigure}[b]{0.48\linewidth}
    \centering
    \zimg{rf_fig_14.png}{width=\linewidth,max height=0.4\textheight,keepaspectratio}{fg_0_9cnzj}{rf_fig_14.png}
    \caption{Impact of changing the CNT diameter in the proposed CNT-based TC-OTAs on (a) DC gain (b) output resistance, (c) average power, and (d) bandwidth}

  \end{subfigure}
  \hfill
  \begin{subfigure}[b]{0.48\linewidth}
    \centering
    \zimg{rf_fig_15.png}{width=\linewidth,max height=0.4\textheight,keepaspectratio}{fg_1_gae9c}{rf_fig_15.png}
    \caption{Impact of changing the CNT diameter in the proposed CNT-based TC-OTAs on (a) DC gain (b) output resistance, (c) average power, and (d) bandwidth}

  \end{subfigure}
  \caption{Impact of changing the CNT diameter in the proposed CNT-based TC-OTAs on (a) DC gain (b) output resistance, (c) average power, and (d) bandwidth}

  \label{fig:group_6dccg}
\end{figure}

\subsection*{Impact of altering the number of CNT pitch (S) on the proposed TC-OTAs' performance}

Inter-CNT pitch (S), like N and DCNT, is an important design consideration for
improving CNTFET performance. The inter-CNT pitch (S) is the distance between
the centers of two adjacent CNTs \cite{ref11,ref12}. It has been noted that if S
is more than 20 nm, the CNTs will carry almost the same currents. The gate
capacitance and drive current of the CNTs in the middle will, however, be lower
than those of the CNTs at the edges due to screening by the nearby CNTs when the
S is smaller. For N = 20 and DCNT = 1.5 nm, the effect of changing CNT pitch on
TC-OTA functionality has been studied. The dc gain, average power, ROUT, and
bandwidth of the suggested CNT-based TC-OTAs are shown to vary with S in Figure
5. While DC gain first increases as S increases, Figure 5(a) shows that it
finally achieves saturation. According to Equation (10), this is explained by an
increase in channel width (W), which first causes an increase in drain current.

\begin{equation}
  W = (N - 1)*S + DCNT
  \label{eq:10}
\end{equation}

However, the inter-CNT screening effect begins to take over as we increase S;
this reduces net gate capacitance, which reduces drain current. But as S is
continuously increased, the inter-CNT screening effect takes over, lowering the
drain current by decreasing the net gate capacitance. When NCNTFET-PMOS-TC-OTA
and PCNTFET-NMOS-TC-OTA are compared, it is evident that the former has a higher
DC gain than the latter. Furthermore, a slight increase in S modifies the DC
gain of the hybrid TC-OTAs. The output resistance of all OTAs decreases as S
increases, as seen in Figure 5(b). Because S grows as channel width increases,
output current also increases, which lowers output resistance
\cite{ref17,ref18,ref19}. As S increases before saturation, the average power
first increases in all three scenarios, as seen in Figure 5(c). Pure CNT-TC-OTA
has the lowest average power because of the above listed issues. The average
power consumption of hybrid OTAs is significantly higher due to the larger
capacitive and short-channel effects in MOSFETs. In all three types of TC-OTAs,
Figure 5(d) shows that bandwidth increases as pitch increases. The increase is
particularly apparent in PCNTFET-NMOS-TC-OTA because of the high drive current
and decreased screening effect brought about by using 12 PCNTFETs rather than
the customary 14 CNTFETs. Compared to pure CNT-TC-OTA and NCNTFET-PMOS-TC-OTA,
the bandwidth enhancement in PCNTFET-NMOS-TC-OTA is less pronounced. This is
justified by the fact that the screening effect is greater because there are
more NCNTFETs \cite{ref35,ref38}.

\begin{figure}[htbp]
  \centering
  \begin{subfigure}[b]{0.48\linewidth}
    \centering
    \zimg{rf_fig_16.png}{width=\linewidth,max height=0.4\textheight,keepaspectratio}{fg_0_uo0s3}{rf_fig_16.png}
  \end{subfigure}
  \hfill
  \begin{subfigure}[b]{0.48\linewidth}
    \centering
    \zimg{rf_fig_17.png}{width=\linewidth,max height=0.4\textheight,keepaspectratio}{fg_1_b98wh}{rf_fig_17.png}
    \caption{Impact of changing the CNT pitch in the proposed CNT-based TC-OTAs on (a) DC gain (b) output resistance, (c) average power, and (d) bandwidth}

  \end{subfigure}
  \caption{Impact of changing the CNT pitch in the proposed CNT-based TC-OTAs on (a) DC gain (b) output resistance, (c) average power, and (d) bandwidth}

  \label{fig:group_wsr60}
\end{figure}

\begin{figure}[htbp]
  \centering
  \begin{subfigure}[b]{0.48\linewidth}
    \centering
    \zimg{rf_fig_18.png}{width=\linewidth,max height=0.4\textheight,keepaspectratio}{fg_0_regyg}{rf_fig_18.png}
    \caption{Impact of changing the CNT pitch in the proposed CNT-based TC-OTAs on (a) DC gain (b) output resistance, (c) average power, and (d) bandwidth}

  \end{subfigure}
  \hfill
  \begin{subfigure}[b]{0.48\linewidth}
    \centering
    \zimg{rf_fig_19.png}{width=\linewidth,max height=0.4\textheight,keepaspectratio}{fg_1_vxeg2}{rf_fig_19.png}
    \caption{Impact of changing the CNT pitch in the proposed CNT-based TC-OTAs on (a) DC gain (b) output resistance, (c) average power, and (d) bandwidth}

  \end{subfigure}
  \caption{Impact of changing the CNT pitch in the proposed CNT-based TC-OTAs on (a) DC gain (b) output resistance, (c) average power, and (d) bandwidth}

  \label{fig:group_2vx1j}
\end{figure}

\subsection*{Impact of altering the gate oxide thickness (TOX) on the proposed TC-OTAs' performance}

The gate oxide functions as a capacitor in a field-effect transistor. The gate
capacitance is lowered by a thicker TOX, which lowers the current drive. Reduced
drain current causes transconductance to drop, which raises output resistance
(ROUT) sa shown in Figure 6(b) \cite{ref40}.
In Figure 6(c), for both the pure and hybrid designs, the average power falls as
TOX rises because the smaller capacitance results in lower drain currents, which
lower power consumption \cite{ref39,ref40,ref41}.
The decrease in transconductance as the oxide thickness rises is the main cause
of the decreasing trend in the bandwidth shown in Figure 6(d) for the pure
CNTFET and hybrid designs. In a CNTFET, a thicker oxide layer decreases the
gate-to-channel capacitance, thereby lowering transconductance by weakening the
gate's control over the channel. As the oxide thickens, the frequency response
deteriorates because the bandwidth of an OTA is directly proportional to
transconductance.On the other hand, the PCNTFET-NMOS-TC-OTA remains flat because
the NMOS load's specific biasing or parasitic capacitances likely limit its
high-frequency performance, making it less susceptible to changes in the gate
dielectric of the CNTFET components \cite{ref6,ref39}.

\begin{figure}[htbp]
  \centering
  \begin{subfigure}[b]{0.48\linewidth}
    \centering
    \zimg{rf_fig_20.png}{width=\linewidth,max height=0.4\textheight,keepaspectratio}{fg_0_qoab4}{rf_fig_20.png}
    \caption{(a) shows the behavior of various operational transconductance amplifier (OTA) topologies' DC Gain as the transistors' oxide...}

  \end{subfigure}
  \hfill
  \begin{subfigure}[b]{0.48\linewidth}
    \centering
    \zimg{rf_fig_21.png}{width=\linewidth,max height=0.4\textheight,keepaspectratio}{fg_1_8dq6k}{rf_fig_21.png}
    \caption{(a) shows the behavior of various operational transconductance amplifier (OTA) topologies' DC Gain as the transistors' oxide...}

  \end{subfigure}
  \caption{(a) shows the behavior of various operational transconductance amplifier (OTA) topologies' DC Gain as the transistors' oxide thickness (TOX) varies between 2 nm and 5 nm. Gate capacitance and drain current are often significantly impacted by the thickness of the gate oxide. Although it might result in leakage, a thinner oxide often improves gate control. The transconductance (gmn) and output resistance (ROCASN) of the transistors, which determine DC gain (AV=gmn* ROUT), seem to be leveling out or being controlled by other dominating elements in the OTA within this particular range, according to the flat lines \cite{ref7,ref39}.}

  \label{fig:group_v6x7o}
\end{figure}

\begin{figure}[htbp]
  \centering
  \begin{subfigure}[b]{0.48\linewidth}
    \centering
    \zimg{rf_fig_22.png}{width=\linewidth,max height=0.4\textheight,keepaspectratio}{fg_0_olutj}{rf_fig_22.png}
    \caption{Impact of changing the oxide thickness in the proposed CNT-based TC-OTAs on (a) DC gain (b) output resistance, (c) average power, and (d) bandwidth}

  \end{subfigure}
  \hfill
  \begin{subfigure}[b]{0.48\linewidth}
    \centering
    \zimg{rf_fig_23.png}{width=\linewidth,max height=0.4\textheight,keepaspectratio}{fg_1_gkvag}{rf_fig_23.png}
    \caption{Impact of changing the oxide thickness in the proposed CNT-based TC-OTAs on (a) DC gain (b) output resistance, (c) average power, and (d) bandwidth}

  \end{subfigure}
  \caption{Impact of changing the oxide thickness in the proposed CNT-based TC-OTAs on (a) DC gain (b) output resistance, (c) average power, and (d) bandwidth}

  \label{fig:group_orhfl}
\end{figure}

\subsection*{Impact of altering the dielectric constant (KOX) on the proposed TC-OTAs' performance}

A thinner gate insulator with a higher dielectric constant may be employed to
further enhance the device's performance. Lately, thin films of high-dielectric
materials such as ZrO2 and LaAlO3 have also been used as insulators in CNFETs.
Table 1 lists dielectric materials commonly considered, along with their
dielectric constants.

\begin{table}[H]
  \centering
  \caption{Given their dielectric constants, the oxides in the following table are taken to be considered for gate insulation.}
  \label{tab:given_their_dielectr}
  \renewcommand{\arraystretch}{1.3}
  \begin{adjustbox}{max width=\linewidth}
    \begin{tabularx}{\linewidth}{|>{\raggedright\arraybackslash}X|>{\raggedright\arraybackslash}X|}
      \hline
      Material & Die-electric value ($\epsilon$r) \\ \hline
      Silicon dioxide (SiO2) & 3.9 \\ \hline
      Hafnium silicate (HfSiO4) & 11 \\ \hline
      Hafnium dioxide (HfO2) & 16 \\ \hline
      Zirconium oxide (ZrO2) & 25 \\ \hline
      Lanthanum aluminate (LaAlO3) & 30 \\ \hline
    \end{tabularx}
  \end{adjustbox}
\end{table}

Increasing gate capacitance increases electrostatic control over the carbon
nanotube channel, which is the main driver of the DC Gain graph's growth and
eventual saturation with respect to the dielectric constant (KOX) in Figure
7(a). The transconductance of the CNTFET devices increases with the dielectric
constant, thereby directly enhancing the voltage gain of the OTA systems.
Nevertheless, this effect follows a diminishing-returns pattern in both the Pure
CNTFET and NCNTFET-PMOS configurations, where the gain saturates beyond a
dielectric constant of roughly 16, indicating that the device has reached a
point at which additional increases in gate coupling no longer significantly
overcome the current output resistance constraints. The PCNTFET-NMOS TC-OTA, on
the other hand, stays noticeably flat, indicating that the typical NMOS load
characteristics, not the gate dielectric characteristics of the P-type CNTFET
components, restrict its gain \cite{ref41,ref42}.
The primary factor leading to the output resistance dropping as the dielectric
constant rises is the inverse relationship between device resistance and gate
coupling in Figure 7(b). By raising gate capacitance, a larger dielectric
constant improves the gate's electrostatic control. This causes the transistor's
drain current to increase significantly and its output resistance to decrease
\cite{ref43,ref44}.
As shown in Figure 7(c), the direct correlation between gate capacitance and
drain current is the main reason for the rise in average power in respect to the
Dielectric Constant. The gate capacitance (COX) increases with the dielectric
constant, improving the carrier density in the carbon nanotube channel and
increasing current flow \cite{ref7}.
The enhancement of the transistor's transconductance, as shown in Figure 7(d),
is the main cause of the bandwidth increase that is visible when the dielectric
constant increases. Greater current drive are enabled by a higher dielectric
constant, which strengthens electrostatic control over the carbon nanotube
channel. The frequency response increases significantly since the bandwidth of
an OTA is exactly proportional to the transconductance of its input transistors
\cite{ref39,ref44,ref45}.

\begin{figure}[htbp]
  \centering
  \begin{subfigure}[b]{0.48\linewidth}
    \centering
    \zimg{rf_fig_24.png}{width=\linewidth,max height=0.4\textheight,keepaspectratio}{fg_0_72u81}{rf_fig_24.png}
    \caption{Impact of changing the dielectric constant in the proposed CNT-based TC-OTAs on (a) DC gain (b) output resistance, (c) average...}

  \end{subfigure}
  \hfill
  \begin{subfigure}[b]{0.48\linewidth}
    \centering
    \zimg{rf_fig_25.png}{width=\linewidth,max height=0.4\textheight,keepaspectratio}{fg_1_srjmf}{rf_fig_25.png}
    \caption{Impact of changing the dielectric constant in the proposed CNT-based TC-OTAs on (a) DC gain (b) output resistance, (c) average...}

  \end{subfigure}
  \caption{Impact of changing the dielectric constant in the proposed CNT-based TC-OTAs on (a) DC gain (b) output resistance, (c) average power, and (d) bandwidth}

  \label{fig:group_l0kdd}
\end{figure}

\begin{figure}[htbp]
  \centering
  \begin{subfigure}[b]{0.48\linewidth}
    \centering
    \zimg{rf_fig_26.png}{width=\linewidth,max height=0.4\textheight,keepaspectratio}{fg_0_4qyre}{rf_fig_26.png}
    \caption{Impact of changing the dielectric constant in the proposed CNT-based TC-OTAs on (a) DC gain (b) output resistance, (c) average...}

  \end{subfigure}
  \hfill
  \begin{subfigure}[b]{0.48\linewidth}
    \centering
    \zimg{rf_fig_27.png}{width=\linewidth,max height=0.4\textheight,keepaspectratio}{fg_1_rfxfz}{rf_fig_27.png}
    \caption{Impact of changing the dielectric constant in the proposed CNT-based TC-OTAs on (a) DC gain (b) output resistance, (c) average...}

  \end{subfigure}
  \caption{Impact of changing the dielectric constant in the proposed CNT-based TC-OTAs on (a) DC gain (b) output resistance, (c) average power, and (d) bandwidth}

  \label{fig:group_nhq1o}
\end{figure}

\subsection*{Stability Evaluation of the suggested TC-OTAs}

Stability assessments of the proposed pure CNT-TC-OTA, hybrid TC-OTAs, and bulk
CMOS-based TC-OTA have been successfully completed by examining their frequency
response in Figure 6. It has been demonstrated that pure CNT-based TC-OTA is
more stable than the others. The findings show that the pure-CNT-TC-OTA has the
maximum gain (75.44 dB) and is exceptionally stable. While the phase margin is
84.5$\circ$ dB. The stability analysis also shows that the hybrid TC-OTAs
(NCNTFET-PMOS-TC-OTA and PCNTFET-NMOS-TC-OTA), which have somewhat smaller
phases and gain margins than the pure CNT-based TC-OTA, are stable. The phase
margin of PCNTFET-NMOS-TC-OTA and NCNTFET-PMOS-TC-OTA circuits have phase
margins of 95.2$\circ$ and 107.6$\circ$, respectively. The phase margin for the
bulk CMOS-TC-OTA is 78$\circ$.

\begin{figure}[htbp]
  \centering
  \begin{subfigure}[b]{0.48\linewidth}
    \centering
    \zimg{rf_fig_28.png}{width=\linewidth,max height=0.4\textheight,keepaspectratio}{fg_0_c4x7v}{rf_fig_28.png}
    \caption{Frequency response of different TC-OTAs.}

  \end{subfigure}
  \hfill
  \begin{subfigure}[b]{0.48\linewidth}
    \centering
    \zimg{rf_fig_29.png}{width=\linewidth,max height=0.4\textheight,keepaspectratio}{fg_1_99ufc}{rf_fig_29.png}
    \caption{Frequency response of different TC-OTAs.}

  \end{subfigure}
  \caption{Frequency response of different TC-OTAs.}

  \label{fig:group_e28g2}
\end{figure}

\begin{figure}[htbp]
  \centering
  \begin{subfigure}[b]{0.48\linewidth}
    \centering
    \zimg{rf_fig_30.png}{width=\linewidth,max height=0.4\textheight,keepaspectratio}{fg_0_5yndk}{rf_fig_30.png}
    \caption{Frequency response of different TC-OTAs.}

  \end{subfigure}
  \hfill
  \begin{subfigure}[b]{0.48\linewidth}
    \centering
    \zimg{rf_fig_31.png}{width=\linewidth,max height=0.4\textheight,keepaspectratio}{fg_1_25lxa}{rf_fig_31.png}
    \caption{Frequency response of different TC-OTAs.}

  \end{subfigure}
  \caption{Frequency response of different TC-OTAs.}

  \label{fig:group_s90gd}
\end{figure}

\subsection*{Parametric Evaluation of the suggested TC-OTAs}

Table 2 displays the comparative and parameterized analysis related to each
TC-OTA that was developed and simulated for this study. The proposed circuits,
which have a VDD of 0.9 Volt and a CL of 0.01 pF, are based on the 32 nm
technology node. Furthermore, Table 1 shows that the pure CNT-based TC-OTA has
the highest gain and the CMOS-based TC-OTA has the lowest. Stability assessments
of the proposed pure CNT-TC-OTA, hybrid TC-OTAs, and bulk CMOS-based TC-OTA have
been successfully completed by examining their frequency response in Figure 9.
It has been demonstrated that pure CNT-based TC-OTA is more stable than the
others. The findings show that the pure-CNT-TC-OTA has the maximum gain (75.44
dB) and is exceptionally stable. While the phase margin is 84.5$\circ$ dB. The
stability analysis also shows that the hybrid TC-OTAs (NCNTFET-PMOS-TC-OTA and
PCNTFET-NMOS-TC-OTA), which have somewhat smaller phases and gain margins than
the pure CNT-based TC-OTA, are stable. The phase margin of PCNTFET-NMOS-TC-OTA
and NCNTFET-PMOS-TC-OTA circuits have phase margins of 95.2$\circ$ and
107.6$\circ$, respectively. The phase margin for the bulk CMOS-TC-OTA is
78$\circ$.

\begin{table}[H]
  \centering
  \caption{CNTFET, hybrid, and bulk TC-OTA performance comparison at CL=0.01 pf, VDD=0.9 V at 32 nm technology node, with N=20, S=20 nm, and DCNT=1.5 nm}
  \label{tab:cntfet__hybrid__and_}
  \renewcommand{\arraystretch}{1.3}
  \begin{adjustbox}{max width=\linewidth}
    \begin{tabularx}{\linewidth}{|c|>{\raggedright\arraybackslash}X|>{\raggedright\arraybackslash}X|>{\raggedright\arraybackslash}X|>{\raggedright\arraybackslash}X|c|}
      \hline
      S. No. & Parameters & PURE CNTFET-TC-OTA & NCNTFET-PMOS-TC-OTA & PCNTFET-NMOS-TC-OTA & CMOS-TC-OTA \\ \hline
      1 & DC gain (dB) & 75.44 & 37.55 & 31.12 & 25.4 \\ \hline
      2 & Bandwidth (GHz) & 3.37 & 3.17 & 4 & 3.6 \\ \hline
      3 & Output resistance (K$\Omega$) & 4.7 & 4.75 & 6.46 & 8 \\ \hline
      4 & Average power (µW) & 38.7 & 43.1 & 57.5 & 62.5 \\ \hline
      5 & Average Slew rate (V/µS) & 1272 & 1207 & 2092 & 2130 \\ \hline
      6 & Phase margin in $\circ$ & 84.5 & 107.6 & 95.2 & 78 \\ \hline
      7 & Gain Margin (dB) & 37.5 & 21 & 32 & 21.1 \\ \hline
      8 & Unity gain frequency (GHz) & 3.03 & 2.88 & 3.62 & 3.2 \\ \hline
      9 & CMRR (dB) & 68.09 & 38.7 & 19.25 & 14.74 \\ \hline
      10 & (FOM)1 & 6559 & 2796 & 2176 & 1445 \\ \hline
      11 & (FOM)2 & 0.657 & 0.558 & 0.626 & 0.54 \\ \hline
    \end{tabularx}
  \end{adjustbox}
\end{table}

\subsection*{Figure of Merits (FOM)1 = [Gain(dB)*UGF(kHZ) / power supply(mV)*power consumption(µW)]}

\subsection*{Figure of Merits (FOM)2 = [Slew rate(V/µS)*CL(pF) / power consumption(µW)]}

\subsection*{Designing Band Pass Filters (BPFs) of the suggested TC-OTAs}

\begin{figure}[htbp]
  \centering
  \begin{subfigure}[b]{0.48\linewidth}
    \centering
    \zimg{rf_fig_32.png}{width=\linewidth,max height=0.4\textheight,keepaspectratio}{fg_0_isc5e}{rf_fig_32.png}
    \caption{(a) Proposed CNT-OTA-Based BPF (b) All four OTAs' frequency responses for the specified BPF}

  \end{subfigure}
  \hfill
  \begin{subfigure}[b]{0.48\linewidth}
    \centering
    \zimg{rf_fig_33.png}{width=\linewidth,max height=0.4\textheight,keepaspectratio}{fg_1_l86zs}{rf_fig_33.png}
    \caption{(a) Proposed CNT-OTA-Based BPF (b) All four OTAs' frequency responses for the specified BPF}

  \end{subfigure}
  \caption{(a) Proposed CNT-OTA-Based BPF (b) All four OTAs' frequency responses for the specified BPF}

  \label{fig:group_ecblm}
\end{figure}

\begin{table}[H]
  \centering
  \caption{Performance comparison of CNTFET, hybrid, and bulk-CMOS-TC-OTA-based BPF at VDD = 0.9 V}
  \label{tab:performance_comparis}
  \renewcommand{\arraystretch}{1.3}
  \begin{adjustbox}{max width=\linewidth}
    \begin{tabularx}{\linewidth}{|c|>{\raggedright\arraybackslash}X|>{\raggedright\arraybackslash}X|>{\raggedright\arraybackslash}X|>{\raggedright\arraybackslash}X|c|}
      \hline
      S. No. & Parameters & PURE CNTFET-TC-OTA & NCNTFET-PMOS-TC-OTA & PCNTFET-NMOS-TC-OTA & CMOS-TC-OTA \\ \hline
      1 & Quality factor (Q) & 0.00257 & 0.0167 & 0.0233 & 0.0348 \\ \hline
      2 & Lower cut-off freq (MHz) & 0.0292 & 0.6 & 0.925 & 1.98 \\ \hline
      3 & Higher cut-off freq (MHz) & 3700 & 1820 & 3040 & 2300 \\ \hline
      4 & Bandwidth (MHz) & 3699.97 & 1819.4 & 3039.07 & 2298.02 \\ \hline
      5 & Central Frequency (MHz) & 9.53 & 30.4 & 71 & 80 \\ \hline
      6 & Capacitance C1 (pF) & 1 & 1 & 1 & 1 \\ \hline
      7 & Capacitance C2 (pF) & 0.01 & 0.01 & 0.01 & 0.01 \\ \hline
      8 & Channel length & 32 nm & 32 nm & 32 nm & 32 nm \\ \hline
    \end{tabularx}
  \end{adjustbox}
\end{table}
